\chapter{Contexte général du projet}
\label{chap:contexte}

\section{Introduction}

Ce chapitre présente le contexte général dans lequel s'inscrit ce stage : l'entreprise
d'accueil, le département Informatique (IT) au sein duquel le stage s'est déroulé, la
problématique à l'origine du projet NetSentry, ainsi que les objectifs poursuivis.

\section{Organisme d'accueil}

\subsection{Présentation de l'entreprise}

Le Mövenpick Hôtel de Tanger fait partie du groupe Mövenpick Hotels \& Resorts, une
chaîne hôtelière internationale opérant dans le secteur de l'hôtellerie et du tourisme.
L'établissement propose des prestations d'hébergement, de restauration et
d'organisation d'événements, destinées principalement à une clientèle d'affaires et de
loisirs.

Le groupe Mövenpick a été fondé en 1948 à Zurich (Suisse) par Ueli Prager, initialement
sous la forme d'un concept de restauration. Son activité hôtelière a débuté en 1973 avec
l'ouverture de ses deux premiers hôtels à Zurich. Depuis septembre 2018, Mövenpick
Hotels \& Resorts est intégralement détenu par le groupe Accor, ce qui lui permet de
bénéficier de son réseau de distribution et de ses outils opérationnels à l'échelle
internationale. Le groupe est aujourd'hui présent dans plusieurs dizaines de pays, avec
une forte présence en Europe et au Moyen-Orient.

L'hôtel de Tanger est structuré en plusieurs départements assurant chacun une fonction
spécifique : la Direction Générale, les Ressources Humaines, la Finance, l'Hébergement
(Front Office, Housekeeping), la Restauration, ainsi que le département Informatique
(IT), au sein duquel s'est déroulé ce stage.

\subsection{Environnement de l'entreprise}

La clientèle du Mövenpick Hôtel de Tanger se compose principalement de voyageurs
d'affaires, de touristes internationaux et de participants à des événements ou
séminaires organisés au sein de l'établissement. L'hôtel fait par ailleurs appel à
différents fournisseurs et prestataires externes, notamment pour la restauration, la
maintenance des équipements, et la gestion des systèmes informatiques métiers (Property
Management System -- PMS), assurée par un prestataire spécialisé externe au département
IT interne. Le marché hôtelier de Tanger est marqué par la présence de plusieurs
établissements internationaux et locaux, notamment dans les segments haut de gamme et
affaires ; le Mövenpick se positionne comme un établissement haut de gamme, misant sur la
qualité de service et son emplacement stratégique dans la ville.

\subsection{Département d'accueil et mission}

Le département Informatique assure la gestion et la maintenance du système
d'information de l'hôtel : réseau informatique, postes de travail, imprimantes et
équipements divers. Les systèmes métiers principaux (PMS, point de vente) sont gérés par
un prestataire externe spécialisé, ce qui recentre l'activité du département IT interne
sur l'infrastructure réseau et le support technique. Il intervient de façon transversale
auprès de l'ensemble des autres départements de l'hôtel et assure la disponibilité et la
sécurité du réseau utilisé par le personnel.

Au sein de ce département, l'étudiant a été chargé de la conception et du développement
d'un outil interne de découverte et de cartographie des équipements connectés au réseau
informatique de l'hôtel, en complément des tâches de support courantes du département.

\section{Contexte}

Un réseau local d'hôtel voit se connecter, au fil du temps, un grand nombre
d'équipements hétérogènes : postes fixes du personnel, imprimantes partagées, caméras de
surveillance, bornes d'accès Wi-Fi, ainsi que des appareils invités (smartphones,
ordinateurs portables) qui se connectent temporairement. Cette diversité et ce
renouvellement permanent rendent difficile le maintien d'un inventaire fiable des
équipements réellement présents sur le réseau à un instant donné.

Avant le déploiement de NetSentry, le département IT du Mövenpick Hôtel de Tanger ne
disposait d'aucun outil de ce type : la connaissance du réseau reposait uniquement sur la
mémoire des techniciens et sur des vérifications ponctuelles et manuelles, sans
historique ni mécanisme d'alerte.

\section{Problématique}

L'absence d'un inventaire automatisé et à jour des équipements connectés pose plusieurs
difficultés concrètes pour le département IT :

\begin{itemize}
    \item Impossibilité de savoir, à un instant donné, combien d'appareils sont
    effectivement actifs sur le réseau et lesquels ;
    \item Absence de traçabilité : aucun historique ne permet de savoir quand un
    équipement est apparu, ni de repérer sa disparition ;
    \item Aucun mécanisme ne permet de détecter et de signaler la connexion d'un
    équipement inconnu ou potentiellement non autorisé ;
    \item Les vérifications, lorsqu'elles ont lieu, sont manuelles, ponctuelles et non
    reproductibles, ce qui les rend peu fiables sur la durée.
\end{itemize}

Ces limitations exposent le réseau interne à un risque de connexions non maîtrisées et
compliquent la tâche du département IT dans le suivi de son propre parc d'équipements. Il
devient donc nécessaire de développer un outil automatisé, léger et autonome, capable de
scanner régulièrement le réseau, de conserver un historique exploitable, et de signaler
toute anomalie de façon proactive.

\section{Objectifs du projet}

Le projet NetSentry vise à répondre à cette problématique à travers les objectifs
suivants :

\begin{itemize}
    \item \textbf{Découvrir} automatiquement les équipements connectés au réseau local
    et les identifier (adresse IP, adresse MAC, fabricant, ports ouverts) ;
    \item \textbf{Centraliser} ces informations dans une base de données persistante,
    consultable via un tableau de bord ;
    \item \textbf{Historiser} les scans afin de suivre l'apparition et la disparition des
    équipements dans le temps ;
    \item \textbf{Détecter et signaler} automatiquement tout équipement détecté pour la
    première fois, afin que l'administrateur IT puisse le vérifier et, le cas échéant, le
    classer comme connu ;
    \item \textbf{Automatiser} l'exécution des scans (planification périodique), sans
    intervention manuelle systématique ;
    \item \textbf{Déployer} la solution de façon conteneurisée (Docker), pour en
    faciliter l'installation et la portabilité sur un serveur de l'hôtel.
\end{itemize}

\section{Conclusion}
\enlargethispage{\baselineskip}

Ce chapitre a posé le cadre général du stage : l'entreprise d'accueil, le département
Informatique, ainsi que la problématique et les objectifs du projet NetSentry. Le
chapitre suivant présente l'étude de l'existant et la méthodologie de travail retenue.
